\item  \textbf{SAT Scores.  [8 pts]} A professor at a small univeristy in New York surveys her introductory statistics students at the beginning of each semester.  One of the questions she asks her students is what was the highest SAT score they received.  The professor has collected data over many semesters and has SAT scores for 353 students.  These data can be found on Blackboard (Assignments - Homework Assignments - Homework 2) in SAT.jmp/SAT.csv/SAT.sas/SAT.xlsx.  Use your favorite software to create a histogram and a boxplot of the SAT scores as well as to calculate summary statistics.  You do not need to give me the output, but you will need it to answer the following questions.
\begin{enumerate}
\item \textbf{ [3 pts]}  Based on the JMP output, describe the overall shape of the distribution of SAT scores for introductory statistics students that attended this small university in New York.\\
\red{The distribution of SAT scores is unimodal, fairly symmetric, with no outliers.}

\item \textbf{  [2 pts]}  Suppose you decide to use a Normal model to describe SAT scores for all introductory statistics students at this small university in New York.  What possible values of the parameters ($\mu$ and $\sigma$) could you use for this Normal model?\\
\red{The best guess we have for the values of the population parameters are the sample statistics.  So, we would use the sample mean, $\bar{y} = 1206.72$ as an estimate for $\mu$.  (Note that since we are dealing with a symmetric, unimodal distribution with no outliers, you could also use the median=1200, but in general you want to estimate the population parameter with the same summary statistic from the sample.)  Similarly, we would use the sample standard deviation, $s = 109.54 $ as an estimate for $\sigma$.}

\item \textbf{ [3 pts]}  The worst score a person can get on the SAT is a 400 and the best score a person can get is a 1600.  A Normal model allows for any value, including values less than 400 and greater than 1600.  Explain why it is still okay to use a Normal model to approximate the distribution of SAT scores. \\
\red{The shape of the histogram (symmetric and unimodal with no outliers) provides evidence that a Normal model would be a good approximation to the distribution of SAT scores.  While the Normal model allows for values from $-\infty$ to $\infty$, we know that almost all values (99.7\%) are within 3 standard deviations of the mean.  Using our estimates for $\mu$ and $\sigma$ from part (b), we know 99.7\% of all observations would be between 878.10 and 1535.34.  400 is well below 878.10 and 1600 is greater than 1535.34, so for this particular Normal model, the probability of observing a value greater than 1600 or less than 400 is so small that the probability is essentially zero.  This is why it is okay to use a Normal model to approximate the distribution of SAT scores.}

\end{enumerate}
